\begin{proposition}[F7: arbitrary coefficients extending over a fixed resolution; Proved]
Let $A,U$ satisfy F5, and fix an SNC resolution $X\to\Spec A$,
with nonempty reduced exceptional curve $D$. There is a constant
$C_X$ such that every finite locally constant $\mathbf F_p$-sheaf
$M$ on $X$ satisfies
\[
 \dim_{\mathbf F_p}H^1(U,M|_U)
       \leq C_X\operatorname{rank}_{\mathbf F_p}M.
\]
The monodromy of $M$ may be nonabelian and may have order divisible
by $p$.
\end{proposition}
\begin{proof}
We give all changes needed to the uniform part of F5, retaining its
notation $V=M\otimes_{\mathbf F_p}\mathcal O_X$, $J$, $Z$, $B_Z$,
and the nearby-cycle filtration. The Artin--Schreier and Frobenius
identities (F5.1)--(F5.2) remain valid: they only use that $M$ is
locally constant on $X$, equivalently that $V$ has its unit
Frobenius structure. We must replace the decomposition into torsion
line bundles by rank-uniform estimates for $V$ itself.

For each smooth exceptional component $C=D_i$, the bundle $V|_C$
is trivialized by a finite \'etale cover of $C$. It is semistable of
degree zero. Indeed its determinant becomes trivial on that cover,
so has degree zero. The pullback of any subbundle becomes a
subbundle of a trivial bundle. Taking its determinant gives a
nonzero map to a finite sum of trivial line bundles, so its degree
is nonpositive. Degree multiplies by the positive covering degree,
so the original subbundle also has nonpositive degree. This proof
works equally for covers of degree divisible by $p$.

It follows that
\[
 H^0(C,V|_C\otimes T)=0\quad\text{if }\deg T<0,
 \qquad h^0(C,V|_C)\leq r,
 \quad r=\operatorname{rank}V.
\]
For the second inequality, evaluation of sections at any closed
point of $C$ is injective: a section vanishing there would give a
line subbundle of positive degree after saturation. Riemann--Roch
therefore gives
\[
 h^1(C,V|_C)\leq r g(C),\qquad
 h^1(C,V|_C\otimes T)=r(g(C)-1-\deg T)\quad(\deg T<0).
 \tag{F7.1}
\]

For the high-pole calculation in F5, replace a layer indexed by
$L_a$ by
\[
 V|_C\otimes(B_Z|_C\otimes P_\nu)\otimes J^{-m}|_C.
\]
The fixed vector bundle in parentheses embeds in a finite sum of
copies of the fixed line bundle $H_\nu$. Semistability and the first
vanishing above show that the displayed layer has no sections as
soon as $\deg H_\nu-m\deg J|_C<0$. Thus the same fixed $m_0$
works, independently of $V$ and $r$.

For the low-pole calculation, the negative-degree filtration of
$J^{-m}/\mathcal O_X$ and $J^{-pm}/J^{-m}$, tensored with $V$,
has no sections. Formula (F7.1) bounds its first cohomology by a
constant times $r$. The semilinear fixed-space lemma in F5 bounds
$H^0(D,R^1\Psi M)$ by the same type of constant. Finally the
normalization sequence for $D$ and (F7.1) give
\[
 h^1(D,V|_D)\leq r\left(\sum_i g(D_i)+\#\{\text{nodes}\}\right).
\]
As in F5, the Artin--Schreier sequence identifies $H^1(D,M)$ with
the Frobenius-fixed subgroup of $H^1(D,V|_D)$, so its
$\mathbf F_p$-dimension is at most the latter $k$-dimension.
The nearby-cycle five-term sequence finishes the proof.
\end{proof}

\begin{lemma}[F7a: a relative pro-$p$ generation criterion; Proved]
Let $1\to P\to G\to\Gamma\to1$ be an exact sequence of profinite
groups, with $P$ pro-$p$ and $\Gamma$ generated by $d$ elements.
Suppose that for every finite simple $\mathbf F_p[G]$-module $M$
there is a uniform bound
$\dim H^1(G,M)\leq C\dim M$. Then $G$ is generated by at most
$\max(d,\lceil C\rceil+1)$ elements. No solvability hypothesis on
$\Gamma$ is needed.
\end{lemma}
\begin{proof}
Put $D_0=\max(d,\lceil C\rceil+1)$. In any finite quotient $F$
of $G$, let $N$ be the image of $P$. Induct on $|N|$ to generate
$F$ by $D_0$ elements. If $N=1$, $F$ is a quotient of $\Gamma$
and the assertion holds. Otherwise choose a minimal nontrivial
$F$-normal subgroup $M\subseteq N$. It is an elementary abelian
$p$-group and a simple $\mathbf F_p[F/M]$-module, by the same
characteristic-subgroup argument as in F4. Induction generates
$F/M$ by $D_0$ elements.

The nonsplit case and the count of nongenerating lifts as
$|Z^1(F/M,M)|$ are exactly those in F4. In the split case the
nonzero $\operatorname{End}_{F/M}(M)$ summand in $H^1(F,M)$
gives the strict bound
\[
 |Z^1(F/M,M)|\leq |M|^{C+1}/p<|M|^{D_0}.
\]
Inflation to $G$ supplies the assumed bound on $H^1(F,M)$.
Thus a generating lift exists. Compactness of $G^{D_0}$, applied
to all finite quotients, completes the argument as in F4.
\end{proof}

\begin{corollary}[F7b: pro-$p$ extensions of quotients from the resolution; Proved]
Let $G$ be a quotient of $\pi_1(U)$ with normal pro-$p$ subgroup
$P$. Suppose the composite
\[
 \pi_1(U)\longrightarrow G\longrightarrow\Gamma=G/P
\]
factors through $\pi_1(X)$. Then $G$ is topologically finitely
generated, with a bound depending only on the fixed resolution $X$.
The group $\Gamma$ need not be abelian, solvable, or prime to $p$.
\end{corollary}
\begin{proof}
Every finite simple $\mathbf F_p[G]$-module is trivial on the normal
pro-$p$ subgroup, by the invariant-vector argument in F6. Its
pullback to $\pi_1(U)$ therefore extends to $X$ under the stated
factorization. Inflation and F7 give the uniform bound required
by F7a.

It remains to bound the generators of $\Gamma$. Proper
specialization over the complete henselian ring $A$ gives
$\pi_1(X)\simeq\pi_1(D)$; passing to the reduced special fibre does
not change finite \'etale covers
\cite[Lemma 58.9.1, tag 0A48]{StacksProperEtaleSpecialization}.
The connected curve $D$ is proper over the algebraically closed
field, so its fundamental group is finitely generated. For a
current stronger primary statement, see
\cite[Theorem 1.2, p.~1]{LaraSrinivasStix2024}, which proves finite
presentation for every connected proper scheme over an
algebraically closed field (and recalls SGA~1 X, Theorem~2.9 for
finite generation). Hence one fixed generator bound works for
all quotients $\Gamma$ of $\pi_1(X)$, and F7a applies.
\end{proof}

This does not assert that the kernel of $\pi_1(U)\to\pi_1(X)$ is
pro-$p$. Tame inertia and more complicated wild quotients may
remain in that kernel. The explicit normal-pro-$p$ hypothesis in
F7b is essential to the argument.
